\subsection{System Instructions and Tool Use at Deployment}
\label{sec:4.2.8}
Training ends and the system's behavior gets set once more. A system instruction is prepended to every request, stating what the deployment is for, what it will decline, and what tone it takes. Tools extend the same arrangement to action: a model can be trained to decide which external calls to make, when to make them, and with what arguments, from a handful of demonstrations per interface \autocite{schick2023toolformer}.

The system instruction is the most editable floor in the entire pipeline, and section~\ref{sec:3.1}'s objection applies to it more literally than to anything else in this chapter. It is a piece of text. It is written by whoever operates the deployment, revised between one request and the next, generally invisible to the person the system is talking to, and binding on that deployment alone. A commitment that lives there is held by the operator, at the operator's pleasure, and the operator is frequently not the party that trained the model. A model tuned with whatever care into whatever refusals is licensed to somebody who then writes the instruction, and what reaches a user is the intersection of what the model will do and what that party asked for. Nothing in the pipeline above survives this stage that the operator wants gone, except to the extent that the training made it expensive to remove.

Tool use changes what is at stake in a refusal. A system that only emits text refuses by declining to say something. A system that can call an interface refuses by declining to do something, and the acts available to it are the acts the deployment wired up — which is how a model reaches the targeting workflow section~\ref{sec:10.3} describes. The floor this book asks for has to hold at the point of action, under an instruction written by the party holding the deployment, on a request that the party has already decided it wants carried out.
